\subsubsection{The lift-adjoint scalar in the selected row metric}
\label{subsec:selected-lift-scalar}

With the kernel norms and terminal evaluation in hand, the remaining normalization is the scalar of a lift followed by its adjoint in the actual row metric.

Fix a lift \(j\) at the reference tuple. In the adelic
calculation write \(F^+\) for the real field, \(E\) for
its CM extension, and \(f=[F^+:\Q]\). The row rank \(r\), definite rank \(n_V\) and stable-range gap satisfy
\[
             r=k+m-1,\qquad n_V=r+t,\qquad t\ge r+5 .
\] The definite source has the
standard self-dual lattice with an integral orthonormal
frame. The local level and Weil vectors use the specified
split half places, with at most one place above each prime
of \(F^+\). The characters intertwine the same local
level actions on the vectors and the forms, and the
Schwartz vector is spherical away from those places.
We use the level, weight, rational Gram and character
identifications of the selected primitive rows. These
conditions apply at the reference tuple, including its
conjugate coefficient realizations.
\label{cor:selected-lift-scalar-1}

\begin{proposition}[Scalar bounds and reference height]\label{prop:selected-reference-height}
For a selected reference sequence \(u\), let \(B_u\) be
the comparison operator in the framed spectral row. With
the common row metric and the period shift
\(\bar\kappa=\overline{\log\kappa}\), its height satisfies
\[
 \operatorname{ht}(B_u)\le h_A(u)+
       \sum_{j=1}^m\bigl(h_B(j)-h_X(j)\bigr)-\bar\kappa .
\]
Here \(h_A\) is the terminal bound, \(h_B(j)\) bounds the
kernel height at lift \(j\), and \(h_X(j)\) bounds the
logarithm of the adjoint scalar \eqref{eq:selected-adelic-scalar}
on the occurring primitive rows. The estimates are uniform on the reference
sequence data, including the coupled end steps.
\end{proposition}

\begin{proof}
\proofstep{Unfold the adjoint in the weighted form and multiplicity metrics.}
Let \(T\) be the full theta lift from the definite form
space with its weighted class metric to the cusp
multiplicity space \(H_W\). This space has the row
metric, which is \(\prod_z d_z\) times the Petersson
multiplicity product, where \(d_z\) is the dimension
of the lowest compact type at \(z\). This metric can
be computed with a unit vector of that type at each
real place by Schur orthogonality. For the finite level group \(P\), the normalized and
unnormalized adjoints on the stated metrics are related by
\[
 T^\dagger=|P|^{-1}T^*,\qquad
 \langle T^\dagger U,T^\dagger U'\rangle
    = |P|^{-2}\langle T^*U,T^*U'\rangle .
\]
For the reference scalar \(|P|=1\).

Let \(U\) lie in a nonzero simultaneous off-level Hecke
part and let \(U'\in H_W\). Let \(C\) be the \(d_C\)-dimensional coefficient module
on the definite side, with orthonormal basis \((e_\nu)\).
Write \(\mathcal R_\nu(U)\) for the scalar reverse
theta integral against \(e_\nu\) in the Petersson
normalization. The source-level transformation law makes
\((\mathcal R_\nu(U))_\nu\) a definite form of the required
character. Put \(D_\infty=\prod_z d_z\). The row metric
contributes \(D_\infty\) to each adjoint, so the weighted
class pairing is
\[
 \langle T^*U,T^*U'\rangle
   =D_\infty^2\sum_{\nu=1}^{d_C}
      \langle\mathcal R_\nu(U),\mathcal R_\nu(U')\rangle_{\rm Pet,def}.
\]
Here the definite Petersson pairing is the sum over its
classes weighted by reciprocal stabilizer orders. The definite quotient is compact,
and cuspidal rapid decay gives the integrability on
the other side.

The convergent Rallis inner-product identity for these
reverse integrals unfolds the product into the finite
oscillator matrix coefficient and the real Fock and
discrete-series coefficients. This is the unitary
doubling identity for cuspidal forms, with the
anisotropic definite member and with the matching
flat section given by the two Schwartz vectors. The
shift is \(s_0=t/2\). In Yamana's Theorem 2, the
doubled skew-Hermitian space has dimension \(2r\),
the definite Hermitian space has dimension \(n_V\),
and its rational Witt index is zero. The range in Yamana's theorem puts the special point in absolute convergence for the Eisenstein series, since
\[
 n_V>2r\quad\Longrightarrow\quad (n_V-r)/2>r/2 .
\] The
Eisenstein series equals the theta integral with
coefficient one for the quotient measure of mass one.

Write \(d\mu_1\) for the definite probability measure and
\(d\mu\) for the definite class-sum measure on the quotient \(Y\).
The conversion in the theta integral is
\[
   \int_Y \theta\,d\mu=\mu(Y)\int_Y\theta\,d\mu_1 .
\] Unfolding and factoring the spherical zeta
integrals gives the partial standard factor at
\(s_0+1/2\). The retained local integral is absolutely
convergent there. The finite zeta evaluation and the two
real evaluations below use the same covolume-one
oscillator measure and real group measures.
\cite{yamana-siegel-weil}\label{cor:selected-lift-scalar-2}

Let \(v_z\) be a unit lowest-type vector and \(H_z^+,H_z^-\) the hook products of the positive and negative partitions of the actual Fock vector at \(z\). Its matrix-coefficient factor is
\[
 J_z=\int |\langle v_z,\pi_z(g)v_z\rangle|^2\,dg .
\] Let \(I_V\) be the index of the source
level in its self-dual stabilizer, \(M_V\) the mass at
that stabilizer, and \(N_\Phi\) the squared norm of
the finite Weil vector.

Write \(S\) for the finite
places omitted in the spherical product, \(D_S\)
and \(D^S\) for the two parts of the normalizing
Euler product, and \(L^S\) for the standard factor
including its unitary character twist. The unfolding
of one reverse integral yields the factor
\[
 I_VM_V\Bigl(\prod_z H_z^+H_z^-J_z\Bigr)
       |P|\,N_\Phi\, L^S/D^S
\]
times the Petersson multiplicity product. The
level factor \(|P|\) counts the classes in the
finite sum that differ by the level group, on which
the two characters cancel. The type determines
\(L^S\) through the spherical eigenvalues of the
first argument. Taking its diagonal entry for a
nonzero \(U\) shows that the resulting scalar is
real and nonnegative. The positive lower estimates
below show it is positive.

In particular the
absolute value of the standard factor equals its
value in that diagonal formula. Let \(A_{\rm rev}\)
denote the displayed factor for one normalized coefficient
basis vector. The reverse products carry \(D_\infty^2\)
from the two adjoints. Since
\(\langle U,U'\rangle_{\rm row}
      =D_\infty\langle U,U'\rangle_{\rm Pet}\),
the sum over the orthonormal coefficient basis is
\[
 \langle T^*U,T^*U'\rangle
     =D_\infty\,d_C A_{\rm rev}\langle U,U'\rangle_{\rm row}.
\]
Here \(A_{\rm rev}\) includes the counted factor \(|P|\). Dividing by \(|P|\)
gives, for every \(U'\),
\(\langle TT^\dagger U,U'\rangle=\lambda\langle U,U'\rangle\),
where
\begin{equation}\label{eq:selected-adelic-scalar}
 \lambda=\Bigl(\prod_z d_z\Bigr)d_C M_V
       \Bigl(\prod_z H_z^+H_z^-J_z\Bigr)
       \frac{|L^S|}{D_SD^S}.
\end{equation}
Here the exact finite-level identity used in the last
line is \(I_V N_\Phi=D_S^{-1}\).
\label{cor:selected-lift-scalar-3}

\proofstep{Bound the finite Euler, level, and stabilizer factors.}

At a split place of residue size \(q\), reducing the
source level gives a parabolic fixing the specified
first \(r\) columns up to a subgroup \(P_q\) of
\(\mathrm{GL}_r(\mathbf F_q)\). Its last \(t\times t\)
block is arbitrary invertible, with \(rt\) free
off-diagonal entries. The relative index and
the squared oscillator norm there are
\[
 \frac{|\mathrm{GL}_{r+t}(\mathbf F_q)|}
      {|\mathrm{GL}_t(\mathbf F_q)|q^{rt}|P_q|},
 \qquad |P_q|q^{-r(r+t)}.
\]
The second equality counts admissible \(r\)-column
matrices, each with a unit-modulus coefficient, in
the integral matrix ball. Of \(r\) columns in dimension \(t+r\), the choices for column
\(j+1\) after \(j\) independent choices are \(q^{t+r}-q^j\).
Their proportion in the full residue matrix space is
\[
 \prod_{j=0}^{r-1}\frac{q^{t+r}-q^j}{q^{t+r}}
       =\prod_{a=1}^r(1-q^{-t-a}),
\]
the reciprocal local normalizer.

At the other places the vector is the
vacuum of squared norm one and the source has its
self-dual level. The split half-place condition allows
multiplication over distinct \(F^+\)-places. This
proves the finite-level identity, also with prescribed
local character factors. The stabilizer index and mass use
the same measure, so the quotient volume at the level
is \(I_VM_V\).

For \(w\) over a spherical place let \(q_w=\mathrm Nw\).
The unramified zeta integral gives an inverse Euler
polynomial whose \(r\) parameters have modulus at most
\(q_w^{(r-1)/2}\) in the normalization before evaluating
at \(s_0+1/2\). This bound uses the spherical part of the
unitarizable cuspidal representation. Consequently
each factor of \(L^S\) is \((1-x_{wa})^{-1}\) with
\[
 |x_{wa}|\le q_w^{(r-1)/2-(1+t)/2}\le q_w^{-3}.
\]
There are at most \(2f\) primes over a rational
prime and \(\sum_{n\ge2}n^{-3}<1\). Thus the Euler
product converges absolutely, and
\(-\log|1-x|\ge-|x|\) gives
\[
 \log|L^S|\ge-\sum_{w\notin S}\sum_{a=1}^r q_w^{-3}
           \ge-2fr\sum_{n\ge2}n^{-3}\ge-2fr .
\]

At a base place of norm \(q\)
the normalizer is
\(\prod_{a=t+1}^{t+r}(1-\epsilon_q^a q^{-a})^{-1}\)
with \(\epsilon_q\) the local quadratic sign in the
doubling normalization. The \(r\) arguments have
absolute value at most \(q^{-2}\). The inequality \(-\log(1-q^{-2})\le2q^{-2}\) and the bound of \(f\) places over each rational prime give, for either positive product \(D\in\{D_S,D^S\}\),
\[
 \log D\le 2fr\sum_p p^{-2}\le4fr .
\] These inequalities are independent of
the finite set omitted and of the occurring Hecke type.

Let \(w_E\) be the number of roots of unity of \(E\).
The identity class alone gives
\(M_V\ge(w_E^{n_V}n_V!)^{-1}\). Indeed an integral
unitary matrix in the standard frame has one
nonzero entry in each column. For a nonzero integral
entry, the mean squared modulus over the complex
embeddings is at least one by the norm and the
arithmetic--geometric mean inequality. Indeed for a column \((a_i)_i\) the relation and positivity give
\[
 1=\sum_i\frac1{[E:\Q]}\sum_\sigma|\sigma(a_i)|^2
   \ge\sum_{i:a_i\ne0}|N_{E/\Q}(a_i)|^{2/[E:\Q]}
   \ge\#\{i:a_i\ne0\}.
\]

The unit length forces exactly one nonzero entry, and equality in
the mean inequality forces all of its conjugates to have the same
absolute value. Thus the one entry has
modulus one at every embedding, so it is a root
of unity. The same argument applies to all columns,
and invertibility forces the positions to form a
permutation. Euler's totient bound \(\varphi(a)^2\ge a/2\), applied to
the roots of unity in \(E\), now gives
\[
              w_E\le 2\varphi(w_E)^2\le 8f^2,\qquad
              M_V\ge\frac{1}{(8f^2)^{n_V}n_V!}.
\]
\label{cor:selected-lift-scalar-4}

\proofstep{Retain the compact hook factors in the lower bound.}

The lowest polynomial is a product of initial minors
in the positive block and reverse initial minors in
the negative block, with exponents successive
differences of the two partitions. The two blocks
use disjoint sets of variables. In covolume-one
Gaussian coordinates the squared Fock norm is
therefore the product of the two hook products.
Each block uses \eqref{eq:selected-fock-matrix}
on its highest matrix coefficient. A determinant
character changes neither its compact normalization
nor the invariant coefficient pairing. This gives
the \(H_z^+H_z^-\) factor in
\eqref{eq:selected-adelic-scalar}.

For the regular half-integral Harish--Chandra
parameter \(\gamma_{z,1}>\cdots>\gamma_{z,r}\)
of the selected discrete series, Schur orthogonality
and the formal-degree formula give
\[
 J_z^{-1}=\frac{1}{V_r}
       \prod_{a<b}\frac{|\gamma_{z,a}-\gamma_{z,b}|}{b-a},
 \qquad
 V_r=\prod_{\ell=1}^r\frac{2\pi^\ell}{(\ell-1)!}.
\]
For compact probability measure the density at the
identity, in the comparison Lie algebra coordinates,
is \(2^{-r(r-1)/2}/V_r\). Indeed the quotient of
\(U(\ell)\) by \(U(\ell-1)\) contributes the unit
sphere of real dimension \(2\ell-1\), with volume
\(2\pi^\ell/(\ell-1)!\) in the column coordinates.
The paired off-diagonal coordinates in the chosen
Lie algebra volume supply the factor \(2^{-r(r-1)/2}\)
in the probability density.

The compact-dual measure
matches that density on \(U(p_z,q_z)\). Its inverse
coefficient integral is the displayed root product
without \(V_r^{-1}\). The density ratio at the identity is
\[
 \frac{2^{-r(r-1)/2}}{2^{-r(r-1)/2}/V_r}=V_r .
\]
Multiplication by this constant converts the compact-dual
coefficient integral to the real measure in the unfolding. The
strict parameter inequalities give integrability and
positivity of \(J_z\). If \(d\) is the degree of the
local Fock polynomial, the parameter width is at
most \(d+n_V+3\), by the lowest-type shift. Each
denominator \(b-a\) is at least one. The volume
formula gives \(V_r\ge r^{-r^2}\), so
\(\log J_z^{-1}\le r^2\log(r(d+n_V+3))\).
\cite{harish-chandra-degree}\label{cor:selected-lift-scalar-5}

Put \(L'=\log(b\,\mathrm{rd}(E)\,g_0g_1 n_V)\).
At the selected sizes,
\(L'\ge5.5\), \(\log n_V\le L'\), and the degree
bound \(d\le3rb\) gives \(\log(d+2)\le2L'\).
Because \(n_V\ge2r\) and \(n_V\ge8\),
\[
 \log\big(r(d+n_V+3)\big)\le\log r+\log n_V+\log(d+2)
    \le4L' .
\]
Summing the real loss over the \(f\) places costs at most
\(4fr^2L'\). The mass loss is at most
\[
 n_V\bigl(4f+\log 2+L'\bigr)
\]
by \(\log w_E\le4f+\log 2\).
The finite loss from the Euler, level and stabilizer factors
of \eqref{eq:selected-adelic-scalar} is at most \(10fr\).
Both dimension factors in \eqref{eq:selected-adelic-scalar}
are positive integers. Adding the losses gives
\begin{equation}\label{eq:selected-hook-lower}
       \log\lambda\ge
        \sum_z\log(H_z^+ H_z^-)-8 f n_V^2 L'.
\end{equation}
In particular the error bound is the same for every
occurring type, unit vector and embedding at a lift.

\proofstep{Descend the scalar at each conjugate embedding of an occurring part.}

At the base embedding the untwisted principal-level
parts span the type decomposition of the full multiplicity
space. They are orthogonal for the row product. The good
type projector \(Q\) is orthogonal and commutes with
this decomposition. The actual lift is \(S_0=QT\).
On a type on which \(TT^*=\lambda I\), the compressed operator
satisfies
\[
 S_0S_0^*=QTT^*Q=\lambda Q .
\]
Thus \eqref{eq:selected-hook-lower} holds on every nonzero
type in the projected image.

At another inducing embedding, use the rational lift
with the selected phase one, its conjugate splitting
characters and its conjugate finite vector. The
conjugate Gram identification gives an isometry of the
definite source forms for the form obtained by applying
that embedding to the rational source pairing. For an
embedding fixing the inducing quadratic field, express the rational form on the definite classes as a sum of geometric theta evaluations with algebraic coefficients.
The same class sum with embedded coefficients evaluates
the new integral, by the conjugate coefficient map and
the covariance of theta evaluation.

Let \({\rm ev}\) denote evaluation on the chosen extremal vector. Equality after this evaluation suffices in the multiplicity space because
\[
 {\rm ev}(v)={\rm ev}(v')\quad\Longrightarrow\quad v-v'\in\ker({\rm ev})=0 .
\] This identifies
the good projection of the new lift with its rational realization up to the unitary change of row coordinates, the action of
a central level element, and a unit-modulus scalar.
The last factor is a product of central character
values, a root of unity, and the unit phase of the
theta comparison. For this unitary map and the central action write their
product as \(U\) and the remaining scalar as \(c\). Relative to
the normalized row Gram matrix \(G\) this gives
\[
 (cU)^*G(cU)=|c|^2 U^*GU=G .
\] At the remaining embeddings
use complex conjugation of both Gram pairings. The
Petersson product is computed anew in the conjugate
model.

For a retained type, take the source vector \(x\)
in the tied rational adjoint image. At the base
embedding its Gram pairing with the rational kernel
is zero. This vanishing is a rational class-sum identity,
so holds at the other embeddings as well. The conjugate
source realization is an isometry, and the rational kernel
spans the Hilbert-model kernel after scalar extension.
Thus the realized source vector is in the orthogonal
complement of the conjugate lift's kernel. The primitive
type and lift identifications give, in its metrics,
\[
 \|S_v x\|_v^2=\lambda_{j,p}(v)\,M_v(x,x).
\]
Here \(M_v\) is the source Gram majorant and the
conjugate realization of the target uses the same
row basis as in the period identity. Both products
carry the reciprocal level index.

In the definite
Hilbert model the majorant is the full coefficient
product divided by that index. The conjugate
target norm is divided by the same index. To obtain
the scalar relation one uses the conjugate
class-sum eigenspace identity on the rational lift and
the conjugate Gram identity on the source. This
identifies the spectral labels before comparing
norms. On \((\ker S_v)^\perp\) the squared singular
values of \(S_v\) are the positive eigenvalues
of \(S_v S_v^*=QTT^*Q\). If \(s_{\min}\) is the smallest positive singular value on the
occurring part, expand a vector \(x\) in its orthonormal right
singular basis. This gives
\[
 \|S_vx\|_v^2=\sum_a s_a^2|x_a|^2
       \ge s_{\min}^2\sum_a|x_a|^2=s_{\min}^2 M_v(x,x).
\] A type present in the
coefficient comparison has some \(x\) with
\(M_v(x,x)>0\), which allows cancellation in the
scalar relation.

Conjugation permutes the labelled weight data in the
hook sum. The negative hooks are positive integers
and may be dropped. The remaining hooks are
\(H_{jz}\) of the fixed positive row partitions.
Consequently for \(C_{\rm sz}\ge8\) a uniform logarithmic lower bound is
\begin{equation}\label{eq:selected-hX}
 h_X(j)=\sum_z\log H_{jz}-C_{\rm sz} f n_V^2 L',
 \qquad \log\lambda_{j,p}(v)\ge h_X(j)
\end{equation}
on all occurring types and at all inducing embeddings.
A type with zero part makes no contribution to the
norm comparison. If a scalar is assigned to that
label in the numerical record, it is assigned the
exponential of the stipulated logarithmic lower bound, as in the
absent-part clause of the scalar record.
\label{cor:selected-lift-scalar-6}

\proofstep{Divide by the lift scalars and add the terminal cost.}

The terminal pairing first uses the weighted class metrics
of Proposition~\ref{prop:terminal-evaluation-height}. At the
same reference level \(P_0\), let \(V\) be the inclusion of
the good row basis and let \(H\) be the Gram matrix of the
principal class sum. With all matrices evaluated at the place in question, the framed norm of a test \(F\) is given by
\[
 \|F\|^2=\frac{\bar F^{t}V^*HV F}{|P_0|}.
\] The principal-level
restriction of the theta evaluation has coordinates \(VF\).
In fact this restriction on full-level forms is the level
idempotent. It cancels the section and cutoff in the
evaluation, since \(V\) already lies in the good level
image. The matrices on both sides are evaluated in the
same basis. The finite-level companion compares the two evaluations on one framed representative of each class. They coincide by equivariance on the remaining points.

The full-level class sum of that evaluation is \(1/|P_0|\)
times its principal-level class sum. For each full-level
orbit, with stabilizer \(S\), its principal-level classes
are the orbits of \(S\) on the finite level quotient
of size \(|P_0|\). Let \(X\) be this finite quotient. Orbit--stabilizer gives
\[
 \sum_{[x]\in S\backslash X}\frac1{|S_x|}
    =\frac1{|S|}\sum_{[x]\in S\backslash X}|Sx|
    =\frac{|P_0|}{|S|}.
\]
The Hermitian squared values are constant on this set
of classes because the level characters have modulus
one. The principal-level sum is \(\|VF\|_H^2\), so
the full-level evaluation has squared norm \(\|F\|^2\).
This gives the norm-nonincreasing evaluation required
in the terminal bound at every infinite place.
\label{cor:selected-lift-scalar-7}

Now use the quotient of this framed test space by the
kernel of its map \(q_0\) to the primitive rows. Its
metric is the norm of the orthogonal representative.
By \eqref{eq:selected-pairing} the terminal pairing
is constant on the fibres of \(q_0\), so it has the
same norm bound in the quotient metric. In the period
comparison the row isomorphism \(\iota\) from this
quotient to the spectral row multiplies its Hermitian
pairing with first argument of type \(p\) by
\[
             \kappa(v)\prod_{j=1}^m\lambda_{j,p}(v).
\]
This is the scalar identity for the same reference
row and word. The other argument is unrestricted.
Thus the type spaces lie in eigenspaces of the
positive self-adjoint comparison of the two metrics.
They span, and the least scalar is a lower bound for
the squared norm ratio on the whole quotient.

Passing to the spectral row dual divides the squared
operator bound by this least scalar. If \(H_j(v)\)
is the logarithmic squared kernel bound including
class sums, the resulting logarithmic squared bound
for a nonzero period map for the sequence \(B_u\) at \(v\) is
\[
 h_A(u)+\sum_{j=1}^m H_j(v)-\log\kappa(v)
                       -\min_p\sum_{j=1}^m\log\lambda_{j,p}(v).
\]
Here the two terminal kernel factors in the coupled
map are counted at their separate lift indices.
Average over the infinite places. The kernel bounds
of Section~\ref{subsec:selected-kernel-jets} and their
CM means give the prices \(h_B(j)\). Using
\eqref{eq:selected-hX} gives
\begin{equation}\label{eq:selected-reference-height}
 \operatorname{ht}(B_u)
    \le h_{\rm num}(u)-\bar\kappa,\qquad
 h_{\rm num}(u)=h_A(u)+\sum_{j=1}^m\bigl(h_B(j)-h_X(j)\bigr).
\end{equation}
The same numerical upper bound may be assigned to a zero period map for the sequence since its operator norm is zero.
The ranks give \(n_V\le N_m\), and \(L'=O(L_a+L)\)
by the selected type-prime and discriminant bounds.

Consequently the allowance in \(h_X(j)\) costs at
most \(O(fN_m^2(L_a+L))\) per lift, uniformly in the
word, place and type. Its hook main term is common
to all sequences on this family. Equation
\eqref{eq:selected-reference-height} uses the exact
numerical prices of the selected comparison after
increasing their one common absolute constant.
\label{cor:selected-lift-scalar-8}
\end{proof}

